% ============================================================
%  CHAPTER 7 — Conclusion and Future Work
%  Logical chain: Contribution summary → Future directions
%  Status: Not started
% ============================================================

\chapter{Conclusion and Future Work}
\label{chap:conclusion}

% ─── 7.1 SUMMARY OF CONTRIBUTIONS ───────────────────────────
\section{Summary of Contributions}
\label{sec:summary}

% Concise synthesis — do NOT introduce new material.
% One paragraph per contribution identified in §1.6.
% Reflect the logical chain: Problem → Gap → Solution → Validation → Contribution.

\lipsum[47]

\begin{criticalinsight}
  % TODO: Final honest self-assessment: does the thesis fully answer the
  % research question (§1.3)? Be precise about what is answered
  % and what is deferred.
\end{criticalinsight}

% ─── 7.2 FUTURE WORK ────────────────────────────────────────
\section{Future Work}
\label{sec:future}

% Should flow directly from §6.5 (Open Problems).
% Structure: near-term (straightforward extensions) vs. long-term (research challenges).

\paragraph{Near-term extensions.}
\begin{itemize}
  \item Hardware validation on the selected platform (removes simulation-only limitation).
  \item Quantitative sim-to-real gap analysis via domain randomization.
  \item Metric threshold grounding from literature-wide meta-analysis (resolves PQ6).
\end{itemize}

\paragraph{Long-term research challenges.}
\begin{itemize}
  \item Generalized contact transition framework that does not require per-platform hand-tuning (§7.2 Methodological Gap).
  \item Long-horizon task sequencing under WBC for service robotics scenarios (§7.4 Application Gap).
  \item Energy-aware WBC formulations for hardware deployment (§7.4 Application Gap).
  \item Tighter integration of learning-based priors into the QP cost structure while preserving constraint guarantees.
\end{itemize}

\lipsum[48]
